\section{SWE-bench: un banco de pruebas para la evaluación de la ingeniería de software}
\subsection{La evolución de Lite y Full}
Graham presenta a continuación SWE-bench: un marco de evaluación en el que un Agente resuelve issues reales de GitHub. Cada muestra incluye el repositorio original, la descripción del issue y la suite de pruebas; la tarea del Agente es comprender la semántica y el contexto, escribir un parche y pasar las pruebas de ese repositorio para verificar su corrección.

SWE-bench se divide en dos categorías: Lite (300 problemas relativamente simples) y Full (2,294 issues de repositorios reales). Lite se concentra en comprender comandos, parches y pruebas básicas; Full extiende el problema a dependencias entre varios archivos, llamadas entre módulos y entornos complejos. El Agente está limitado a una ventana de contexto, por lo que necesita diseñar estrategias eficaces de prompt, recuperación por índice y razonamiento local.

\subsection{Proceso de evaluación y métricas}
\begin{figure}[H]
\centering
\includegraphics[width=\textwidth]{frame-02-swebench.jpg}
\caption{Panel de estadísticas de evaluación de SWE-bench, que muestra la tasa de éxito de los conjuntos de datos Lite y Full. \protect\footnotemark}
\end{figure}
\footnotetext{Intervalo de tiempo en el video: 00:15:20--00:15:25.}

El proceso de evaluación consta de tres pasos: 1) el Agente lee el issue y recupera el contexto relevante en el repositorio; 2) genera y aplica un parche (edita archivos, ejecuta \texttt{git diff}); 3) ejecuta las pruebas objetivo y archiva los resultados junto con los registros de salida. El sistema calcula la tasa de éxito según \texttt{pass/fail} y retroalimenta las muestras fallidas al proceso de entrenamiento para mejorar el prompt y las estrategias de herramientas.

\begin{table}[H]
\centering
\begin{tabular}{p{0.35\textwidth} p{0.58\textwidth}}
\toprule
Dimensión & Descripción \\
\midrule
Tipo de tarea & Extraer el objetivo del issue, ubicar los archivos relevantes, generar un parche y ejecutar las pruebas correspondientes. \\
\midrule
Lite & 300 puntos propensos a error, orientados a verificar la estructura del prompt y la capacidad de corrección rápida; los mejores Agentes actuales alcanzan una tasa de éxito de entre 40\%--50\%. \\
\midrule
Full & 2,294 desafíos, incluidas llamadas asíncronas, migración de dependencias y parches de seguridad; la tasa de éxito general actual es inferior a 20\%. \\
\bottomrule
\end{tabular}
\caption{Dimensiones de diseño de SWE-bench y resultados actuales.}
\end{table}

\begin{knowledgebox}{Énfasis de la evaluación de SWE-bench}
SWE-bench enfatiza el ciclo completo de ingeniería: el Agente debe entregar el patch, ejecutar las pruebas y proporcionar los registros; no basta con dar el diff. La división entre Lite y Full permite a los equipos de investigación optimizar primero las capacidades básicas antes de extenderse a escenarios complejos.
\end{knowledgebox}

\begin{importantbox}{Indicadores cuantitativos}
Graham enfatiza: una tasa de éxito de 40\% en el nivel Lite ya significa que el Agente puede sustituir parte de las tareas repetitivas, pero la baja tasa de éxito en el nivel Full indica que todavía estamos en una era de «semiautomatización».
\end{importantbox}

\begin{warningbox}{Límites de cobertura}
SWE-bench evalúa principalmente la codificación y las pruebas; no incluye tareas más amplias de ingeniería de software como el diseño, la revisión de código, el despliegue y el mantenimiento, de modo que las capacidades evaluadas siguen inclinadas hacia problemas de tipo corrección.
\end{warningbox}

\subsection{Resumen de la sección}
\begin{itemize}
    \item SWE-bench ofrece un espectro doble que va de lo simple a lo complejo, y le da al Agente un objetivo que puede alcanzarse de manera gradual.
    \item Los indicadores actuales siguen muy por debajo de los de un ingeniero humano, y todavía se necesitan nuevas evaluaciones para etapas más amplias del ciclo de vida del software.
\end{itemize}
